\documentclass[11pt]{article}

\usepackage[T1]{fontenc}
\usepackage{lmodern}
\usepackage{amsmath,amssymb,amsthm}
\usepackage{booktabs}
\usepackage[hidelinks,pdftitle={Complete Tight-Set Classification for Deletion Speed 12},pdfauthor={Xiaojian Zeng},pdfsubject={Lonely Runner tight-set classification},pdfkeywords={Lonely Runner Conjecture, tight instance, exact rational certificate}]{hyperref}
\usepackage[margin=1in]{geometry}

\newtheorem{theorem}{Theorem}
\newtheorem{lemma}{Lemma}
\newcommand{\dist}[1]{\left\lVert #1\right\rVert}
\newcommand{\LR}{\operatorname{LR}}

\title{Complete Tight-Set Classification for Deletion Speed \(12\)}
\author{Xiaojian Zeng\thanks{PBC School of Finance, Tsinghua University,
Beijing 100190, China. Email: \texttt{xiaojianzeng@foxmail.com}.}}
\date{}

\begin{document}
\maketitle

\begin{abstract}
For integers \(n\geq 13\) and \(w>n-1\), put
\[
 V_{n,w}=([n-1]\setminus\{12\})\cup\{w\},\qquad
 \LR(V)=\max_{0\leq t<1}\min_{v\in V}\dist{vt},
\]
where \(\dist{x}\) denotes distance to the nearest integer.  We prove that
\(\LR(V_{n,w})=1/n\) if and only if \((n,w)=(14,24)\).  Published
divisibility, effective-bound, and census results reduce the infinite
classification to twenty pairs.  Exact rational interval-cover certificates
then determine the value for every residual pair; a separate complete
breakpoint enumeration gives the same values.  The result resolves the fixed
deletion-speed-\(12\) specialization, not the broader middle-range
classification for arbitrary deletion speeds.
\end{abstract}

\section{Statement and context}

The Lonely Runner Conjecture originated in work of Wills
\cite{Wills1967}.  In the integer-speed normalization used here, a set of
\(n-1\) distinct positive speeds is called \emph{tight} when its lonely-runner
value is exactly \(1/n\).  Goddyn and Wong classified the tight sets obtained
from the baseline \([n-1]=\{1,\ldots,n-1\}\) by replacing one speed with an
integer multiple of itself \cite{GoddynWong2006}.

Zhang made the complementary finiteness result effective, gave an exact
census through \(n=45\) in a bounded speed range, and asked for the remaining
middle-range classification \cite{Zhang2026}.  The census contains the tight
pair \((14,24)\), but does not settle fixed deletion speed \(12\) beyond its
range.  The contribution here is to combine Zhang's universal reduction with
an exact treatment of every case left for this one deletion speed.  No
priority claim beyond this comparison is made.

\begin{theorem}\label{thm:main}
For every pair of integers \(n\geq13\) and \(w>n-1\),
\[
  \LR\bigl(([n-1]\setminus\{12\})\cup\{w\}\bigr)=\frac1n
  \quad\Longleftrightarrow\quad (n,w)=(14,24).
\]
Moreover, \(\LR(V_{14,24})=1/14\), while the exact lonely-runner value
of each of the twenty residual pairs in Table~\ref{tab:values} is one of
\(2/59,2/61,2/65,2/67\).
\end{theorem}

\section{Universal reduction}

We first reduce the quantified statement to a finite list.  This step is
essential: the subsequent exact computation alone would not prove
Theorem~\ref{thm:main}.

Put \(s=n-12\).  Suppose first that \(13\leq n\leq24\).  Since
\(2\cdot12>n-1\), Zhang's divisibility proposition forces \(12\mid w\)
\cite[Proposition~1.3]{Zhang2026}; write \(w=12m\).  Here \(m\geq2\).
The Goddyn--Wong criterion \cite[Theorem~2.3]{GoddynWong2006}, restated as
Zhang's Theorem~1.1 \cite{Zhang2026}, says that tightness is equivalent,
outside exceptions which do not occur here, to
\[
 \gcd(12,b)>1\qquad(s\leq b\leq ms-1)
 \tag{1}\label{eq:gcdcriterion}
\]
For \(n=13\), the interval in \eqref{eq:gcdcriterion} contains \(1\).
For \(n=14,m=2\), it is \(\{2,3\}\), so the criterion holds.  For
\(n=14,m\geq3\), it contains \(5\).  Finally, for
\(n=15,\ldots,24\), the interval already contains, respectively,
\[
 5,5,5,7,7,11,11,11,11,13,
\]
at \(m=2\), and it continues to contain that integer for larger \(m\).
Each displayed obstruction is coprime to \(12\).  Thus this whole range
contributes exactly \((14,24)\), and the same criterion proves
\(\LR(V_{14,24})=1/14\).

Now let \(n\geq25\).  For a residue \(u\), define
\[
 F_n(u)=\max\{i\leq n-1:i\equiv u\pmod {12}\},\qquad
 I(n)=\min_{u\in\{1,5,7,11\}}F_n(u).
\]
The deleted value \(12\) is irrelevant to these four unit residue classes.
Zhang's effective bound \cite[Theorem~1.6]{Zhang2026} implies that a tight
pair must satisfy
\[
 n\leq72,
 \qquad
 w\leq B(n):=\frac{48I(n)}{2(n-12)-I(n)}.
 \tag{2}\label{eq:bound}
\]
For \(25\leq n\leq45\), direct substitution gives
\(\lfloor B(n)\rfloor\leq74<10n\).  Zhang's exact census
\cite[Theorem~1.8(i) and Table~1]{Zhang2026} therefore covers every possible
\(w\) in this range and contains no deletion-speed-\(12\) pair.

For \(46\leq n\leq72\), substitution in \eqref{eq:bound} yields the
following nonempty ranges; at \(n=49\) and every \(52\leq n\leq72\), one
has \(\lfloor B(n)\rfloor<n\).
\[
\begin{array}{c@{\qquad}c@{\qquad}c@{\qquad}c}
\toprule
n&I(n)&\lfloor B(n)\rfloor&\text{possible }w\\
\midrule
46&35&50&46\text{--}50\\
47&35&48&47\text{--}48\\
48&37&50&48\text{--}50\\
50&41&56&50\text{--}56\\
51&41&53&51\text{--}53\\
\bottomrule
\end{array}
\tag{3}\label{tab:reduction}
\]
Consequently, \eqref{tab:reduction} is the complete twenty-pair residual list.  Values
\(n>72\) are excluded directly by \eqref{eq:bound}.

\section{Exact certification of the residual values}

The upper certificates use the following elementary description.  It also
gives a short verifier independent of breakpoint enumeration.

\begin{lemma}[Interval-cover certificate]\label{lem:cover}
Let \(V\) be a finite set of positive integers and \(0\leq q\leq1/2\).  For
each \(v\in V\),
\[
 \{t\in[0,1]:\dist{vt}\leq q\}
 =\bigcup_{k=0}^{v}
 \left[\frac{k-q}{v},\frac{k+q}{v}\right]\cap[0,1].
 \tag{4}\label{eq:intervals}
\]
If labeled intervals from the right-hand sides of \eqref{eq:intervals} form
a chain covering \([0,1]\), then \(\LR(V)\leq q\).  If additionally a
rational \(t_0\) satisfies
\(\min_{v\in V}\dist{vt_0}=q\), then \(\LR(V)=q\).
\end{lemma}

\begin{proof}
The nearest integer to \(vt\) is some \(k\in\{0,\ldots,v\}\), which proves
\eqref{eq:intervals}.  A cover says that for every \(t\) at least one speed
has distance at most \(q\), proving the upper bound.  Evaluation at \(t_0\)
gives the matching lower bound.
\end{proof}

For every row of Table~\ref{tab:values}, the ancillary file
\texttt{evidence/lr\_interval\_cover\_results.json} stores a labeled chain
as in Lemma~\ref{lem:cover} and all distances at the stated \(t_0\).  There
are \(6{,}305\) intervals and \(964\) lower-witness distances in total.  The
verifier checks with rational arithmetic that the first interval begins at
zero, every next interval meets the current frontier, the last reaches one,
each label generates its stored endpoints by \eqref{eq:intervals}, and the
minimum witness distance is \(q\).

\begin{table}[ht]
\centering
\small
\caption{The complete exact finite evaluation.  Pairs in a row share the
displayed value and attaining time.}\label{tab:values}
\begin{tabular}{c c p{0.60\textwidth}}
\toprule
\(\LR(V_{n,w})\)&attaining \(t_0\)&pairs \((n,w)\)\\
\midrule
\(1/14\)&\(1/14\)&\((14,24)\)\\
\(2/59\)&\(5/59\)&\((46,46),(46,48),(46,49),(46,50),(47,48)\)\\
\(2/61\)&\(5/61\)&\((46,47),(47,47),(48,48),(48,50)\)\\
\(2/65\)&\(27/65\)&\((48,49),(50,50),(50,51),(50,52),(50,54),\)
\newline \((50,55),(50,56),(51,51),(51,52)\)\\
\(2/67\)&\(28/67\)&\((50,53),(51,53)\)\\
\bottomrule
\end{tabular}
\end{table}

A structurally separate program applies Zhang's breakpoint lemma
\cite[Lemma~2.2]{Zhang2026}.  It enumerates, for all \(v<w\) in \(V\),
\[
 \frac{2a+1}{2v},\qquad \frac{b}{v+w},\qquad \frac{b}{w-v}
\]
over the complete index ranges in that lemma, evaluates the minimum speed
distance at every distinct candidate, and maximizes exactly.  It obtains the
same value and least attaining time in all twenty-one cases.  This program
does not establish the universal reduction; conversely, the interval-cover
program neither imports nor reads the breakpoint implementation or its
results.

The proof-carrying separation between an exact serialized certificate and a
verifier which reconstructs its finite objects is also used in other recent
finite classifications; compare the graph-theoretic certificates of
\cite[Section~6]{ZengEtAl2026}.  Here the upper object is instead the explicit
rational cover \eqref{eq:intervals}.

For every residual row of Table~\ref{tab:values}, the displayed value is
strictly greater than \(1/n\).  Thus none of the twenty residual pairs is
tight.  Together with the universal reduction and the positive
Goddyn--Wong criterion at \((14,24)\), this proves
Theorem~\ref{thm:main}.

\section*{Reproducibility and assistance disclosure}

The complete certificates and standard-library Python verifiers accompany
the manuscript.  From the project root, the two decisive computations are
reproduced by
\begin{verbatim}
python3 evidence/lr_breakpoints.py
python3 evidence/lr_interval_cover.py
\end{verbatim}
Both use \texttt{fractions.Fraction}; no floating-point tolerance or random
seed is used.  The interval-cover output contains every interval and every
lower-witness distance, so its mathematical checks do not rely on a search
cache.  The breakpoint output supplies an independent exact cross-check.

OpenAI Codex assisted with literature search, source-scope reconstruction,
development of the exact programs, adversarial audit organization, and
initial manuscript preparation.  Language-model output is not used as
mathematical evidence: the conclusion rests on the cited theorems, the
displayed reduction, and the released exact certificates and verifiers.  No
proof assistant was used.

\bibliographystyle{plain}
\bibliography{references}

\end{document}
